\subsection{Locally compact topological vector spaces}

THEOREM 3. --- Let K be a complete non-discrete valued division ring. If E is a Hausdorff topological vector space over K, which is such that some neighbourhood V of 0 in E is precompact, then E is of finite dimension. If \(E \neq \{0\}\) , then both K and E are locally compact.

In proving the first assertion, we need consider only the case when E is complete; for E is an everywhere dense subspace of its completion \(\hat{E}\) , and the closure \(\overline{V}\) of V in \(\hat{E}\) is compact and is a neighbourhood of 0 in \(\hat{E}\) (GT, III, § 3.4, prop. 7).

We can suppose then that there is a compact neighbourhood V of 0 in E. Let \(\alpha \in \mathbf{K}\) be such that \(0 < |\alpha| < 1\) ; then there are finitely many points \(a_i \in \mathbf{V}\) such that

\[
\mathrm{V} \subset \bigcup_ {i} (a _ {i} + \alpha \mathrm{V}).
\]

Let M be the finite dimensional subspace of E generated by the \(a_{i}\) ; it is closed in E (I, p.~14, cor. 1). In the Hausdorff topological vector space E/M the canonical image of V is a compact neighbourhood W of 0, such that \(W \subset \alpha W\) ; hence \(\alpha^{-1} W \subset W\) , and, by induction on n, \(\alpha^{-n} W \subset W\) for every positive integer n.~As W is absorbent, we conclude that W = E/M; and thus E/M is compact. To complete the proof of the first assertion in the theorem, it is sufficient, therefore, to establish the following lemma.

Lemma 1. --- Any compact topological vector space E over a non-discrete valued division ring, is just the set \(\{0\}\) .

Since E is complete we can suppose K is complete (I, p.~6). If \(E \neq \{0\}\) then E contains a line that is closed in E (I, p.~14, cor. 1) and therefore compact. This line is isomorphic to \(K_{s}\) (I, p.~12, prop. 2) and hence K must be compact. Now the mapping \(\xi \mapsto |\xi|\) of K in R is continuous and thus the image of K must be bounded, on the other hand there exists \(\gamma \in K\) with \(|\gamma| > 1\) , and the set \(|\gamma^{n}| = |\gamma|^{n}\) , \(n \in N\) , is unbounded. This contradiction shows that \(E = \{0\}\) .

To prove the second assertion in the theorem, if \(E \neq \{0\}\) then from the first part of the theorem E is isomorphic to \(K_{s}^{n}\) with n \textgreater{} 0; now K is complete, hence so is E, and thus E is locally compact. But \(K_{s}\) is isomorphic to a line in E (I, p.~12, prop. 2) which is necessarily closed in E (I, p.~14, cor. 1); it follows that K is locally compact.

Remark. --- The result of th. 3 is no longer true if K is a discrete division ring as is shown by the example of R (with the usual topology) considered as a topological vector space over the discrete field Q.

